\documentclass[UTF8,fontset=fandol]{ctexart}

\usepackage[margin=2cm]{geometry}
\usepackage{amsmath,amssymb,mathtools}
\usepackage{booktabs}
\usepackage{tabularx}
\usepackage{enumitem}

\newcommand{\R}{\mathbb{R}}
\newcommand{\rank}{\operatorname{rank}}
\newcommand{\col}{\operatorname{col}}
\newcommand{\Range}{\operatorname{Range}}
\newcommand{\diag}{\operatorname{diag}}

\title{Cai et al.\ (2017) Lemma 3 中 $\bar{x}$ 的有界性与收敛性}
\author{}
\date{}

\begin{document}

\maketitle

\begin{quote}
来源：He Cai, Frank L. Lewis, Guoqiang Hu, Jie Huang,
\emph{The adaptive distributed observer approach to the cooperative output
regulation of linear multi-agent systems}, Automatica 75 (2017), 299--305。

本文对应原文 Lemma 3 的证明，重点解释方程 (16)--(19)，以及原文中
\emph{``Then, $\bar{x}$ is bounded by Remark 2''} 这一步为什么需要补充说明。
\end{quote}

\section{结论先行}

\begin{enumerate}[leftmargin=2em]
  \item $\smash{P^{T}A^{T}AP}$ 一般\textbf{不是 Hurwitz}。由 SVD 分解，
  \[
  P^{T}A^{T}AP
  =
  \begin{bmatrix}
    \bar{A}^{T}\bar{A} & 0\\
    0 & 0_{n-k}
  \end{bmatrix},
  \]
  其中 $\bar{A}^{T}\bar{A}>0$，但后 $n-k$ 个特征值为零。

  \item 因此，不能直接对整个 $\bar{x}$ 使用论文 Lemma 1 或 Remark 2；
  这些结论要求对应的标称矩阵 Hurwitz。

  \item $\bar{x}$ 的有界性应通过值域与零空间的分块分析得到：
  值域分量受到 $-\varepsilon\bar{A}^{T}\bar{A}$ 的指数稳定作用；
  零空间分量的极限系统是积分器，但其驱动项指数衰减，因此收敛到某个常数。

  \item 只有在先证明 $\bar{x}$ 有界之后，才能由
  \[
  d(t)=\varepsilon P^{T}(A^{T}A-\mathcal{A}^{T}\mathcal{A})P\bar{x}
  +\varepsilon P^{T}\tilde{A}^{T}b
  \]
  推出 $d(t)$ 指数趋于零。若反过来用 $d(t)\to0$ 证明 $\bar{x}$ 有界，
  就会形成循环论证。
\end{enumerate}

\section{Lemma 3 的设定}

设
\[
A\in\R^{m\times n},\qquad b\in\R^{m},
\]
并满足
\[
\rank(A)=\rank(A,b)=k.
\]
因此线性方程 $Ax=b$ 一致可解，但当 $k<n$ 时，解一般不唯一。

令 $\mathcal{A}(t)$ 为有界、分段连续的时变矩阵，并定义
\[
\tilde{A}(t)\triangleq\mathcal{A}(t)-A.
\]
假设 $\tilde{A}(t)$ 以速率 $\alpha>0$ 指数趋于零，即存在常数 $c_A>0$ 使
\[
\|\tilde{A}(t)\|
\leq c_Ae^{-\alpha(t-t_0)},\qquad t\geq t_0.
\]

考虑系统
\begin{equation}
\dot{x}
=-\varepsilon\mathcal{A}(t)^{T}
\bigl(\mathcal{A}(t)x-b\bigr),
\qquad \varepsilon>0.
\tag{12}
\end{equation}

目标是证明：对任意初值，解在 $[t_0,\infty)$ 上有界，并且存在某个
$x^*$ 满足
\[
Ax^*=b,\qquad x(t)-x^*\to0
\]
指数收敛。这里的 $x^*$ 可以依赖于初值；当 $A$ 不满列秩时，不能预先指定唯一的解。

\section{SVD 分解与 $P^{T}A^{T}AP$ 的结构}

取正交矩阵 $P=[P_1\;P_0]\in\R^{n\times n}$，使得
\[
AP=
\begin{bmatrix}
\bar{A} & 0_{m\times(n-k)}
\end{bmatrix},
\qquad
\bar{A}\in\R^{m\times k}.
\]

由于 $\rank(A)=k$，矩阵 $\bar{A}$ 列满秩，因此
\[
G\triangleq\bar{A}^{T}\bar{A}>0.
\]
于是
\begin{equation}
P^{T}A^{T}AP=
\begin{bmatrix}
G&0\\
0&0_{n-k}
\end{bmatrix}.
\tag{13}
\end{equation}

同时，由 $AP_0=0$ 可得
\begin{equation}
P^{T}A^{T}b=
\begin{bmatrix}
\bar{A}^{T}b\\
0
\end{bmatrix}.
\tag{14}
\end{equation}

所以，$P^{T}A^{T}AP$ 只是半正定矩阵，而不是 Hurwitz 矩阵。只有在 $k=n$
时，它才没有零特征值。

由一致性条件，
\[
b\in\Range(A)=\Range(\bar{A}).
\]
因此存在唯一的
\[
\bar{x}_1^*=G^{-1}\bar{A}^{T}b
\]
满足
\[
\bar{A}\bar{x}_1^*=b.
\]
令任意 $\bar{x}_2^*\in\R^{n-k}$，并定义
\[
x^*=P
\begin{bmatrix}
\bar{x}_1^*\\
\bar{x}_2^*
\end{bmatrix}.
\]
则 $Ax^*=b$。

\section{方程 (16) 的推导}

令
\[
\bar{x}=P^{T}x.
\]
将系统 (12) 左乘 $P^{T}$，并在右端加一项、减一项
$\varepsilon P^{T}A^{T}Ax$，得到
\begin{equation}
\begin{aligned}
\dot{\bar{x}}
&=-\varepsilon P^{T}\mathcal{A}^{T}\mathcal{A}x
  +\varepsilon P^{T}\mathcal{A}^{T}b\\
&=-\varepsilon P^{T}A^{T}AP\bar{x}
  +\varepsilon P^{T}A^{T}b+d(t),
\end{aligned}
\tag{16}
\end{equation}
其中
\begin{equation}
d(t)
=\varepsilon P^{T}(A^{T}A-\mathcal{A}^{T}\mathcal{A})P\bar{x}
 +\varepsilon P^{T}\tilde{A}^{T}b.
\tag{D}
\end{equation}

注意
\[
\begin{aligned}
A^{T}A-\mathcal{A}^{T}\mathcal{A}
&=A^{T}A-(A+\tilde{A})^{T}(A+\tilde{A})\\
&=-A^{T}\tilde{A}-\tilde{A}^{T}A-\tilde{A}^{T}\tilde{A}.
\end{aligned}
\]
因此，$A^{T}A-\mathcal{A}^{T}\mathcal{A}$ 以速率 $\alpha$ 指数趋于零。

但是，式 (D) 中还含有 $\bar{x}(t)$。所以目前只能说：
\begin{itemize}
  \item 若已知 $\bar{x}$ 有界，则 $d(t)$ 有界并以速率 $\alpha$ 指数趋于零；
  \item 不能先假定 $d(t)\to0$，再用它证明 $\bar{x}$ 有界。
\end{itemize}

\section{为什么不能直接套用 Remark 2}

论文 Remark 2 的结论适用于形如
\[
\dot{z}=\varepsilon Fz+F_2(t)
\]
且 $F$ Hurwitz 的系统。

在这里，标称矩阵为
\[
-P^{T}A^{T}AP
=-
\begin{bmatrix}
G&0\\
0&0
\end{bmatrix},
\]
它含有零特征值。因此不能把整个 $\bar{x}$ 当作 Remark 2 中的 $z$ 直接处理。

论文中的
\begin{quote}
\emph{Then, $\bar{x}$ is bounded by Remark 2}
\end{quote}
省略了零空间分量的有界性证明。严格的证明应先对值域和零空间进行分块，
或者等价地对误差 $x-x^*$ 进行投影分解。

\section{非循环的有界性证明}

\subsection{关键：不要把 $d(t)$ 当作外部输入}

取上面构造的任意解 $x^*$，满足 $Ax^*=b$，并令
\[
z=P^{T}(x-x^*)
=\col(z_1,z_2).
\]
由于
\[
\mathcal{A}(t)x^*-b
=(A+\tilde{A}(t))x^*-b
=\tilde{A}(t)x^*,
\]
可得
\begin{equation}
\dot{z}=-\varepsilon C(t)z+r(t),
\tag{C}
\end{equation}
其中
\[
C(t)\triangleq P^{T}\mathcal{A}(t)^{T}\mathcal{A}(t)P,
\qquad
r(t)\triangleq-\varepsilon P^{T}\mathcal{A}(t)^{T}\tilde{A}(t)x^*.
\]

\begin{equation}
\begin{aligned}
C(t)&\to
C_\infty\triangleq
\begin{bmatrix}
G&0\\
0&0
\end{bmatrix},\\
\|C(t)-C_\infty\|+\|r(t)\|
&=O(e^{-\alpha(t-t_0)}).
\end{aligned}
\tag{C1}
\end{equation}

下面对 (C1) 的衰减估计给出详细推导。

\subsubsection*{估计 $\|C(t)-C_\infty\|$}

将 $\mathcal{A}(t)=A+\tilde{A}(t)$ 代入 $C(t)$：

\[
\begin{aligned}
C(t) &= P^T(A+\tilde{A})^T(A+\tilde{A})P \\
&= P^T(A^TA + A^T\tilde{A} + \tilde{A}^TA + \tilde{A}^T\tilde{A})P
\end{aligned}
\]

减去 $C_\infty = P^TA^TAP$：

\[
C(t)-C_\infty = P^T(A^T\tilde{A} + \tilde{A}^TA + \tilde{A}^T\tilde{A})P
\]

取范数，利用 $P$ 正交（$\|P\|=\|P^T\|=1$）和 $\|XY\|\leq\|X\|\|Y\|$：

\[
\begin{aligned}
\|C(t)-C_\infty\|
&\leq \|A^T\tilde{A}\| + \|\tilde{A}^TA\| + \|\tilde{A}^T\tilde{A}\| \\
&\leq 2\|A\|\,\|\tilde{A}(t)\| + \|\tilde{A}(t)\|^2
\end{aligned}
\]

代入 $\|\tilde{A}(t)\|\leq c_A e^{-\alpha(t-t_0)}$，并注意 $e^{-2\alpha\tau}\leq e^{-\alpha\tau}$（$\tau\geq0$）：

\[
\|C(t)-C_\infty\|
\leq 2\|A\|c_A e^{-\alpha(t-t_0)} + c_A^2 e^{-2\alpha(t-t_0)}
\leq (2\|A\|c_A + c_A^2)\, e^{-\alpha(t-t_0)}
\]

\subsubsection*{估计 $\|r(t)\|$}

\[
r(t) = -\varepsilon P^T\mathcal{A}(t)^T\tilde{A}(t)x^*
\]

取范数：

\[
\begin{aligned}
\|r(t)\|
&\leq \varepsilon \|P^T\|\,\|\mathcal{A}(t)^T\|\,\|\tilde{A}(t)\|\,\|x^*\| \\
&= \varepsilon \|\mathcal{A}(t)\|\,\|\tilde{A}(t)\|\,\|x^*\|
\end{aligned}
\]

由于 $\|\mathcal{A}(t)\| = \|A+\tilde{A}(t)\| \leq \|A\| + \|\tilde{A}(t)\| \leq \|A\| + c_A$（常数，与 $t$ 无关），代入得：

\[
\|r(t)\| \leq \varepsilon(\|A\|+c_A)\|x^*\|\,c_A e^{-\alpha(t-t_0)}
\]

\subsubsection*{合并}

两项均以 $O(e^{-\alpha(t-t_0)})$ 衰减，故存在常数 $c_0>0$（例如 $c_0 = 2\|A\|c_A + c_A^2 + \varepsilon(\|A\|+c_A)\|x^*\|c_A$）使得

\[
\|C(t)-C_\infty\| + \|r(t)\| \leq c_0 e^{-\alpha(t-t_0)}
\]

由此，令 $\Delta(t) = -\varepsilon(C(t)-C_\infty)$，则 $\|\Delta(t)\| \leq \varepsilon c_0 e^{-\alpha(t-t_0)}$，与 $\|r(t)\|$ 一起满足 (C3)。

定义
\[
A_0\triangleq-\varepsilon C_\infty,
\]
则 (C) 可以改写为
\begin{equation}
\dot{z}=A_0z+\Delta(t)z+r(t),
\tag{C2}
\end{equation}
且由 (C1) 存在常数 $c_0>0$ 使得
\[
\|\Delta(t)\|+\|r(t)\|
\leq c_0e^{-\alpha(t-t_0)}. \tag{C3}
\]

这里的关键是：$A_0$ 虽然不是 Hurwitz，但它是半稳定的，且零特征值是半单的。
事实上，
\[
e^{A_0t}
=
\begin{bmatrix}
e^{-\varepsilon Gt}&0\\
0&I_{n-k}
\end{bmatrix},
\qquad
\|e^{A_0t}\|\leq1,\quad t\geq0. \tag{C4}
\]

\subsection{用变分公式和 Gronwall 先证明有界性}

由 (C2) 的变分常数公式，
\[
z(t)
=e^{A_0(t-t_0)}z(t_0)
+\int_{t_0}^{t}e^{A_0(t-s)}
\bigl(\Delta(s)z(s)+r(s)\bigr)\,ds.
\]
结合 (C3)--(C4)，得到
\[
\|z(t)\|
\leq\|z(t_0)\|
+c_0\int_{t_0}^{t}e^{-\alpha(s-t_0)}\|z(s)\|\,ds
+c_0\int_{t_0}^{t}e^{-\alpha(s-t_0)}\,ds.
\]

下面明确说明 Gronwall 不等式在此处的使用。积分形式的 Gronwall
不等式是指：若非负函数 $y(t)$ 满足
\[
y(t)\leq a+\int_{t_0}^{t}b(s)y(s)\,ds,
\qquad a\geq0,\quad b(s)\geq0,
\]
则
\[
y(t)\leq a\exp\left(\int_{t_0}^{t}b(s)\,ds\right).
\]
在当前问题中，先估计上式中的最后一项：
\[
c_0\int_{t_0}^{t}e^{-\alpha(s-t_0)}\,ds
\leq \frac{c_0}{\alpha}.
\]
因此，令
\[
y(t)\triangleq\|z(t)\|,\qquad
a\triangleq\|z(t_0)\|+\frac{c_0}{\alpha},\qquad
b(s)\triangleq c_0e^{-\alpha(s-t_0)},
\]
就得到
\[
y(t)\leq a+\int_{t_0}^{t}b(s)y(s)\,ds.
\]
应用上述 Gronwall 不等式，并利用
\[
\int_{t_0}^{t}b(s)\,ds
=c_0\int_{t_0}^{t}e^{-\alpha(s-t_0)}\,ds
\leq\frac{c_0}{\alpha},
\]
可得显式估计
\begin{equation}
\|z(t)\|
\leq
\left(\|z(t_0)\|+\frac{c_0}{\alpha}\right)
e^{c_0/\alpha},
\qquad t\geq t_0.
\tag{C5}
\end{equation}
右端与 $t$ 无关，故
\begin{equation}
\sup_{t\geq t_0}\|z(t)\|<\infty.
\tag{B}
\end{equation}
于是 $x(t)=x^*+Pz(t)$ 也有界。

这一步完全没有使用原文定义的 $d(t)$，因此不存在“先假设 $\bar{x}$ 有界”
再证明 $d(t)$ 衰减的循环。

\subsection{再证明收敛速度}

由 (B)，式 (D) 中的 $\bar{x}$ 有界，因此
\[
\|d(t)\|\leq c_de^{-\alpha(t-t_0)}.
\]
于是 (16) 的分块形式为
\begin{equation}
\dot{\bar{x}}_1
=-\varepsilon G\bar{x}_1+\varepsilon\bar{A}^{T}b+d_1(t),
\tag{17a}
\end{equation}
\begin{equation}
\dot{\bar{x}}_2=d_2(t).
\tag{17b}
\end{equation}

对第一式令
\[
\tilde{x}_1=\bar{x}_1-\bar{x}_1^*.
\]
因为 $-G$ Hurwitz，
\[
\dot{\tilde{x}}_1=-\varepsilon G\tilde{x}_1+d_1(t),
\]
从而 $\tilde{x}_1(t)\to0$ 指数收敛。对任意 $\varepsilon>0$，其收敛速率为某个正数；
若
\begin{equation}
\varepsilon>
\frac{\alpha}{\lambda_{\min}(G)},
\tag{R}
\end{equation}
则可以使值域分量至少以速率 $\alpha$ 收敛。

再看零空间分量。由 $d_2(t)=O(e^{-\alpha t})$，
\[
\bar{x}_2(t)
=\bar{x}_2(t_0)+\int_{t_0}^t d_2(\tau)\,d\tau
\]
收敛到某个有限常数 $\bar{x}_2^*$，且
\[
\bar{x}_2(t)-\bar{x}_2^*=O(e^{-\alpha t}).
\]

因此取
\[
x^*=P
\begin{bmatrix}
\bar{x}_1^*\\
\bar{x}_2^*
\end{bmatrix},
\]
就有
\[
Ax^*=b,\qquad
x(t)-x^*
=P
\begin{bmatrix}
\bar{x}_1(t)-\bar{x}_1^*\\
\bar{x}_2(t)-\bar{x}_2^*
\end{bmatrix}
\to0
\]
指数收敛。这就是原文式 (18)--(19) 的结论。

\section{两种分块视角的关系}

直接对 $\bar{x}=P^{T}x$ 分块，可以得到
\[
\dot{\bar{x}}_1
=-\varepsilon G\bar{x}_1+\varepsilon\bar{A}^{T}b+d_1(t),
\qquad
\dot{\bar{x}}_2=d_2(t).
\]
这有助于说明：
\begin{itemize}
  \item $\bar{x}_1$ 是由正定矩阵 $G$ 稳定的值域分量；
  \item $\bar{x}_2$ 属于 $\ker A$，标称系统中没有恢复力；
  \item $\bar{x}_2$ 是否有界，取决于 $d_2$ 是否可积。
\end{itemize}

但由于 $d_2$ 的定义中含有 $\bar{x}$，这组方程单独使用时会暴露出循环。
以 $x-x^*$ 为中心化变量后，先用 (C2)--(C3) 得到有界性，再推出 $d(t)$
指数衰减，逻辑更加完整。

\section{与 Lemma 4 的衔接}

Lemma 4 中令
\[
Q_i(t)
=S_i(t)^{T}\otimes
\begin{bmatrix}
I&0\\
0&0
\end{bmatrix}
-I_q\otimes
\begin{bmatrix}
A_i&B_i\\
C_i&D_i
\end{bmatrix},
\]
以及
\[
b_i=\operatorname{vec}
\begin{bmatrix}
E_i\\
F_i
\end{bmatrix}.
\]

由 Lemma 2，$\tilde{S}_i(t)=S_i(t)-S$ 指数趋于零，因此
\[
Q_i(t)\to Q_i
\]
也指数趋于零。Assumption 3 保证极限方程
\[
Q_i\chi_i^*=b_i
\]
可解。于是可以将 Lemma 3 应用于
\[
\mathcal{A}(t)=Q_i(t),\qquad A=Q_i,\qquad b=b_i.
\]
得到在线求解器 (22) 的解 $\zeta_i(t)$ 有界，并收敛到某个满足
$Q_i\chi_i^*=b_i$ 的解。若增益 $\mu_3$ 足够大，则收敛速率至少可以达到
$\tilde{S}_i(t)$ 的指数衰减速率。

\section{关键提醒}

\begin{itemize}
  \item $P^{T}A^{T}AP$ 一般不是 Hurwitz，而是
  $\diag(\bar{A}^{T}\bar{A},0)$。
  \item 不能直接对整个 $\bar{x}$ 使用 Lemma 1 或 Remark 2。
  \item $d(t)\to0$ 需要以 $\bar{x}$ 有界为前提，不能用它反过来证明 $\bar{x}$ 有界。
  \item 零空间分量不是由 Hurwitz 稳定性控制，而是因为其驱动项可积，最终收敛到常数。
  \item 当 $A$ 不满列秩时，极限解 $x^*$ 通常不唯一，并且会依赖于初始条件。
  \item 若只要求指数收敛，任意 $\varepsilon>0$ 足够；若要求收敛速率至少为
  $\alpha$，需要满足类似 (R) 的条件。
\end{itemize}

\appendix

\section{矩阵指数、变分公式与 Gronwall 不等式}

\subsection{矩阵指数}

\subsubsection*{定义}

对任意方阵 $A \in \R^{n\times n}$，\textbf{矩阵指数}定义为幂级数

\[
e^{A} = \sum_{k=0}^{\infty} \frac{A^k}{k!},
\qquad
e^{At} = \sum_{k=0}^{\infty} \frac{(At)^k}{k!}.
\]

该级数对任意 $t$ 绝对收敛，且 $e^{At}$ 关于 $t$ 光滑。

\subsubsection*{基本性质}

\begin{enumerate}
  \item $e^{A \cdot 0} = I$（单位阵）；
  \item $\displaystyle \frac{d}{dt} e^{At} = A e^{At} = e^{At} A$；
  \item $e^{A(t+s)} = e^{At} e^{As}$（半群性质）；
  \item 若 $A$ 与 $B$ 可交换（$AB=BA$），则 $e^{A+B} = e^A e^B$。
\end{enumerate}

由性质 2，齐次方程 $\dot{x} = Ax$ 的解为 $x(t) = e^{A(t-t_0)} x(t_0)$。

\subsubsection*{特征值的关系}

矩阵指数最重要的性质之一是：\textbf{$A$ 的特征值决定了 $e^{At}$ 的特征值}。

\begin{enumerate}
  \item 若 $\lambda$ 是 $A$ 的特征值，则 $e^{\lambda t}$ 是 $e^{At}$ 的特征值（重数相同）；
  \item 对应的特征向量不变：若 $A v = \lambda v$，则 $e^{At} v = e^{\lambda t} v$。
\end{enumerate}

\subsubsection*{推导（可对角化情形）}

若 $A$ 可对角化：$A = P \Lambda P^{-1}$，$\Lambda = \operatorname{diag}(\lambda_1,\dots,\lambda_n)$，则

\[
e^{At} = P e^{\Lambda t} P^{-1}, \qquad
e^{\Lambda t} = \operatorname{diag}(e^{\lambda_1 t}, \dots, e^{\lambda_n t}).
\]

$e^{At}$ 相似于 $e^{\Lambda t}$，故特征值就是 $e^{\lambda_i t}$。对不可对角化的情形，通过 Jordan 标准型可得相同结论。

\subsubsection*{几个直接推论}

\begin{itemize}
  \item $e^{At}$ 总是可逆的，且 $(e^{At})^{-1} = e^{-At}$。这是因为 $\det(e^{At}) = e^{\operatorname{tr}(A)t} \neq 0$。
  \item 若 $A$ 的特征值实部全为负（Hurwitz），则 $e^{At} \to 0$（当 $t \to \infty$）。
  \item 若 $A$ 有零特征值，则 $e^{At}$ 对应特征值为 $1$——不衰减。这正是本文 $A_0$ 的情况：$-\varepsilon G$ 的特征值为负（衰减），$0$ 矩阵的特征值为 $0$（对应 $e^{0\cdot t}=1$，不衰减）。
  \item 若 $A$ 对称，则 $e^{At}$ 也对称，特征值均为实数，且 $e^{At}$ 正定（当 $t$ 有限时）。
\end{itemize}

\subsubsection*{分块对角阵的情形}

若 $A = \operatorname{diag}(A_1, A_2, \dots, A_m)$ 为分块对角，则

\[
e^{At} = \operatorname{diag}\bigl(e^{A_1 t}, e^{A_2 t}, \dots, e^{A_m t}\bigr).
\]

\subsubsection*{本文的关键例子}

在有界性证明中，

\[
A_0 = -\varepsilon C_\infty
= -\varepsilon \begin{bmatrix} G & 0 \\ 0 & 0 \end{bmatrix}
= \operatorname{diag}(-\varepsilon G,\; 0_{n-k}),
\]

其中 $G = \bar{A}^T\bar{A} > 0$。由分块对角性质：

\[
e^{A_0 t} = \begin{bmatrix} e^{-\varepsilon G t} & 0 \\ 0 & e^{0 \cdot t} \end{bmatrix}
= \begin{bmatrix} e^{-\varepsilon G t} & 0 \\ 0 & I_{n-k} \end{bmatrix}.
\]

由于 $G$ 正定，$-\varepsilon G$ 的特征值全为负，故 $e^{-\varepsilon G t} \to 0$（指数衰减）；而零空间部分的 $e^{0\cdot t} = I$ 保持恒等。

\subsubsection*{为什么 $\|e^{A_0 t}\| \leq 1$}

这里使用的是\textbf{谱范数}（即算子 2-范数，记作 $\|\cdot\|_2$）。对任意对称矩阵，谱范数等于其谱半径：

\[
\|M\|_2 = \max_{i} |\lambda_i(M)|.
\]

$A_0$ 是对称的（因为 $G$ 对称），故 $e^{A_0 t}$ 也对称。其特征值为：

\begin{itemize}
  \item $e^{-\varepsilon \lambda_i(G) t}$，其中 $\lambda_i(G) > 0$ 是 $G$ 的特征值。由于 $t \geq 0$，有 $0 < e^{-\varepsilon \lambda_i(G) t} \leq 1$；
  \item $1$（重数 $n-k$），来自零空间分量 $I_{n-k}$。
\end{itemize}

因此所有特征值的绝对值均不超过 $1$，谱半径为 $1$，即

\[
\|e^{A_0 t}\|_2 = 1, \qquad t \geq 0.
\]

特别地，$\|e^{A_0 t}\|_2 \leq 1$ 成立。需注意：若使用其他范数（如 Frobenius 范数），该结论不一定成立；但本文的 Gronwall 论证只需算子范数满足该有界性即可。

即 $A_0$ 是\textbf{半稳定}的（值域分量衰减，零空间分量有界但不衰减）。这正是不能直接套用 Remark 2（要求 Hurwitz）的原因，也是整个证明需要分块处理的根源。

\subsection{变分公式（常数变易法）}

考虑非齐次线性常微分方程

\[
\dot{x}(t) = A x(t) + f(t), \qquad x(t_0) = x_0,
\]

其中 $A \in \R^{n\times n}$ 为常矩阵，$f(t)$ 分段连续。其解可表示为

\begin{equation}
x(t) = e^{A(t-t_0)} x_0 + \int_{t_0}^{t} e^{A(t-s)} f(s)\,ds.
\tag{V}
\end{equation}

这就是\textbf{变分公式}（variation of constants formula，亦称 Duhamel 原理或常数变易法）。

\subsubsection*{推导思路}

令 $x(t) = e^{A(t-t_0)} y(t)$，则 $y(t_0) = x_0$。求导：

\[
\dot{x} = A e^{A(t-t_0)} y + e^{A(t-t_0)} \dot{y}
= A x + e^{A(t-t_0)} \dot{y}.
\]

与 $\dot{x} = Ax + f(t)$ 对比得 $e^{A(t-t_0)} \dot{y} = f(t)$，即 $\dot{y} = e^{-A(t-t_0)} f(t)$。积分：

\[
y(t) = y(t_0) + \int_{t_0}^{t} e^{-A(s-t_0)} f(s)\,ds
= x_0 + \int_{t_0}^{t} e^{-A(s-t_0)} f(s)\,ds.
\]

代回 $x(t) = e^{A(t-t_0)} y(t)$ 即得 (V)。

\subsubsection*{在有界性证明中的应用}

本文在 (C2) 中处理的是

\[
\dot{z} = A_0 z + \Delta(t) z + r(t),
\]

其中 $A_0$ 为常矩阵，而 $\Delta(t)z + r(t)$ 扮演了 $f(t)$ 的角色。应用变分公式：

\[
z(t) = e^{A_0(t-t_0)} z(t_0)
+ \int_{t_0}^{t} e^{A_0(t-s)} \bigl(\Delta(s) z(s) + r(s)\bigr)\,ds.
\]

这正是文中 Gronwall 论证的起点。

\subsection{Gronwall 不等式}

\subsubsection*{积分形式}

若非负函数 $y(t)$ 满足

\[
y(t) \leq a + \int_{t_0}^{t} b(s) y(s)\,ds,
\qquad a \geq 0,\; b(s) \geq 0,
\]

则

\[
y(t) \leq a \exp\!\left(\int_{t_0}^{t} b(s)\,ds\right).
\]

\subsubsection*{微分形式（等价）}

若 $y(t) \geq 0$ 且 $\dot{y}(t) \leq \beta(t) y(t)$，则

\[
y(t) \leq y(t_0) \exp\!\left(\int_{t_0}^{t} \beta(s)\,ds\right).
\]

\subsubsection*{在本文中的使用}

取 $y(t) = \|z(t)\|$，$a = \|z(t_0)\| + c_0/\alpha$，$b(s) = c_0 e^{-\alpha(s-t_0)}$。由于

\[
\int_{t_0}^{t} b(s)\,ds
= c_0 \int_{t_0}^{t} e^{-\alpha(s-t_0)}\,ds
\leq \frac{c_0}{\alpha} < \infty,
\]

Gronwall 给出与 $t$ 无关的一致有界估计：

\[
\|z(t)\| \leq \left(\|z(t_0)\| + \frac{c_0}{\alpha}\right) e^{c_0/\alpha}.
\]

这正是 (C5) 和 (B) 的来源。

\subsection{小结}

变分公式 + Gronwall 不等式是处理"标称矩阵非 Hurwitz 但半稳定 + 扰动指数衰减"这类问题的标准工具组合：

\begin{enumerate}
  \item 用变分公式将解表示为"自由演化 + 扰动卷积"；
  \item 自由演化部分由 $e^{A_0 t}$ 的范数控制（此处 $\|e^{A_0 t}\| \leq 1$）；
  \item 扰动卷积部分通过 Gronwall 吸收进指数因子，得到一致有界性。
\end{enumerate}

\end{document}
